\section{계획 오차와 시간대 이동의 보조 분석}
\label{sec:legacy}
이 절의 실험은 \textbf{기존 Laplace BAT}와 튜닝 LightGBM으로 수행한 선행 보강 분석이다. 새 조건부 정규 BAT로 다시 실행하지 않았으므로 앞 절의 피크 개선을 시간대 이동의 개선으로 연결하지 않는다. 또한 실제 공장에서 계획을 변경한 결과가 아니라, 입력 변화와 모델 기반 반사실 계산의 민감도를 살펴보는 분석이다.

\subsection{생산계획 입력 오차}
학습 자료와 실제 전력 목표를 고정한 채 예측일의 계획 입력만 바꾸었다. 생산량 오차는
\begin{equation}
q^{\mathrm{input}}_{d,h}=q_{d,h}\exp(\sigma Z_{d,h}-\sigma^2/2),
\quad Z_{d,h}\sim\mathcal N(0,1)
\end{equation}
로 두어 곱셈 오차의 평균이 1이 되게 했다. $\sigma=0.05,0.1,0.2$마다 날짜별 10회 반복하고, 반복 점수는 날짜 안에서 먼저 평균하였다. 0인 시각은 그대로 유지된다. 시간 오차는 전체 계획을 한 시간 앞이나 뒤로 옮기되 날짜 밖으로 나가는 생산을 00시 또는 23시에 남겨 총량을 보존했다. 이는 계획 입력의 오기입 검사이며 변경 후 실제 전력을 생성한 실험이 아니다.

\begin{table}[t]
\centering\small
\begin{tabular}{@{}lrr@{}}
\toprule
입력 조건 & 기존 BAT MAE & 튜닝 M2 MAE \\
\midrule
오차 없음 & 7.820 & 11.048 \\
생산량 $\sigma=0.05$ & 7.820 & 11.049 \\
생산량 $\sigma=0.10$ & 7.821 & 11.049 \\
생산량 $\sigma=0.20$ & 7.823 & 11.056 \\
1시간 이른 계획 & 8.264 & 11.915 \\
1시간 늦은 계획 & 7.878 & 11.741 \\
\bottomrule
\end{tabular}
\caption{기존 모형의 계획 입력 오차 민감도(kW). 동일 53일의 5,016개 관측 슬롯을 사용하며 반복으로 평가 표본 수를 늘리지 않는다.}
\label{tab:plan-error}
\end{table}

기존 BAT의 MAE는 생산량 $\sigma=0.2$에서 0.03\%, 한 시간 이른 계획에서는 5.68\% 증가했다. 로그 채널과 유형별 기본곡선을 함께 사용하는 모델이 생산량의 작은 변화보다 가동 시각 변화에 민감하다는 설명과 부합한다. 그러나 계획 취소·설비 고장·제품 전환은 가동 여부 자체를 바꿀 수 있으므로 이 결과로 그러한 상황의 강건성을 주장하지 않는다.

\subsection{시간대 이동과 정오 제약}
후보 계획은 하루 생산 총량을 보존하고 날짜 경계를 넘지 않는다. 학습기간의 유형별 시작 시각 빈도에 Dirichlet 평활화를 적용하고, 주간·야간의 양수 생산량을 각각 Gamma 적률로 근사했다. 블록 이동은 새로운 시작 시각의 사후예측확률이 5\% 이상인 경우에만, 인접 시간 이동은 더 낮은 요금 시간대이며 이동 후 생산량이 해당 구간의 Gamma 95분위 이내인 경우에만 허용했다. 이 조건들은 과거 관측 지지 범위를 이용한 제약이며 설비 능력과 작업 순서를 직접 측정한 제약은 아니다.

이 보조 시뮬레이션의 실행 오차는 $q^{\mathrm{exec}}_h=q'_h\exp(\sigma Z_h)$를 사용했다. 입력 오기입 실험의 평균 1 보정과 달리 중앙값이 1인 로그정규 오차이며, 실현 생산 총량은 표본마다 달라질 수 있다. 원계획과 후보계획의 오차는 같은 분포에서 독립적으로 생성한다. 후보 식별자에 난수를 고정하여 정오 제약으로 다른 후보를 제거해도 남은 후보의 통계는 유지한다.

사후 모수와 실행 오차 표본에서 계산한 평균구조 곡선을 $f_t^{(s)}(q)$라 하고, 중간·최대부하 시간대를 $\mathcal D$라 하자. 기존 구현의 권고 조건은
\begin{equation}
\frac1S\sum_{s=1}^S\mathbf1\!\left[
 \max_{t\in\mathcal D}f_t^{(s)}(q')
 <\max_{t\in\mathcal D}f_t^{(s)}(q)\right]\ge0.95
\label{eq:decision}
\end{equation}
이다. 이 곡선에는 식~\eqref{eq:path}의 새로운 날짜 효과와 AR 잔차를 추가하지 않는다. 따라서 식~\eqref{eq:decision}은 완전한 미래 관측 경로에 대한 성공확률이 아니라, 기존 평균구조와 실행 오차 아래의 모형 기반 판단이다. 여러 후보가 통과하면 전력량요금과 수요요금 일할 대용치 변화의 표본 평균이 가장 작은 계획을 선택했다.

정오 제약은 $q'_{12}\le q_{12}$로 적용하여 12시의 추가 생산만 금지했다. 정오 완전 비가동을 강제한 것은 아니다. 후보가 없는 1일을 무조치로 포함하여 허용·금지 조건 모두 가동일 39일을 분모로 사용했다. 이 39일은 실제 피크 평가가 가능한 37개 가동일과 목적 및 마스크가 다르다.

\begin{table}[t]
\centering\footnotesize
\setlength{\tabcolsep}{3pt}
\begin{tabular}{@{}lrr@{}}
\toprule
$\sigma=0.1$ & 정오 증가 허용 & 증가 금지 \\
\midrule
대상 / 권고일 & 39 / 20 & 39 / 19 \\
권고일 피크 감소 중앙값 (kW) & 3.38 & 2.11 \\
권고일 전력량요금 중앙값 (원) & $-4{,}433$ & $-4{,}399$ \\
전체일 전력량요금 평균 (원/일) & $-2{,}150$ & $-1{,}891$ \\
\bottomrule
\end{tabular}
\caption{기존 BAT의 정오 제약 민감도. 음의 비용 변화는 모델상 감소를 뜻한다. 권고하지 않은 날은 변화 0으로 집계하며 실제 청구액 절감이 아니다.}
\label{tab:noon}
\end{table}

정오 제약 이후 권고일의 피크 감소는 작아졌지만 전체 가동일 평균 전력량요금 감소의 약 88\%가 남았다. 선택된 계획은 10시에서 09시 이동 12일, 16시에서 17시 이동 6일, 블록 전체를 한 시간 앞당기는 이동 1일이었다. 같은 날 생산 총량의 보존이나 요금 시간대 회피만으로 현장 실행 가능성이 보장되지는 않는다. 실제 기본요금은 월 최대수요 및 과거 피크의 래칫 바닥과 관계되므로 수요요금 일할 대용치와 구분해야 한다. 신규 조건부 모형에서는 이 의사결정 계산을 다시 검증해야 한다.
